\scanpage{147}\mspage{2}
\={m} of $2$-categories everything seems so clear -- ``one only has to follow the tip of one's nose''\dots) and finally the absence of a definition of a functor going from (semi) simplicial sets to $n$-groupoids. This functor corresponds to a geometric intuition so clear that a theory which did not include it would seem like a hoax to me. Perhaps by \inserted{trying to }writ\deleted{\unsure{ing}}\unsure{e} (As\unsure{:|.} a sort of Christmas wish list !) in a fairly complete and explicit way the \deleted{relations} notions \deleted{\unsure{that one}}which one would like to have at one's disposal and the relations (functors, equivalences and so on) which must link them, one would finally arrive at a sort of sufficiently developed axiomatic description \deleted{\unsure{to}}which, either \unsure{gives|give} the key to an explicit ad-hoc construction, or \inserted{which allows} one at a pinch to state and prove an existence and uniqueness theorem for a theory of the desired type.

Furthermore, not having understood Segal's idea in your last letter (which I generously sent to Illusie\dots), I do not see how you define Picard $n$-categories -- but no matter. As for ``strict'' Picard $n$-categories, all I ask of them is that they ultimately form an \inserted{$(n+1)$-}category $(n+1)$-equivalent to that of chain complexes of length\unsure{,} $n$. Agreed ? I thank you
% archive sheet number: 147
